\chapter{Spectral Graph Theory}

\begin{center}\itshape
Can one hear the shape of a drum?\\
\upshape(Mark Kac, 1966)
\end{center}

In 1845 a student at the University of K\"onigsberg, Gustav Kirchhoff, aged twenty-one, stated the two laws that every
electrician knows: the currents into a junction add up to zero, and the voltages around a closed loop add up to zero. Two
years later he solved the problem they pose. Given a network of wires with resistances and a battery, how much current
flows in each wire? The answer is a quotient of two determinants, and Kirchhoff noticed something strange about these
determinants: when every resistance is $1$, they count the \emph{spanning trees} of the network, the ways of choosing
wires that connect every junction without forming a loop. This was the first theorem of spectral graph theory, in the city
whose bridges had started graph theory a century earlier.

The second life of the subject began with telephones. In the 1970s engineers wanted networks in which every group of
subscribers has many connections to the rest, so that no small set of cables could isolate a large set of users, while
each switch has only a few cables. Such graphs, the \emph{expanders}, were shown to exist by counting (Pinsker, 1973), and an
explicit family was found the same year by Grigory Margulis using group theory. In the 1980s Noga Alon, Vitali Milman and
others found the key: a graph is a good expander exactly when the second eigenvalue of its matrix is far from the first.
Expanders are now everywhere in computer science, in error-correcting codes, in pseudorandom generators, in the
proof that one can find one's way out of an undirected maze with logarithmic memory (Reingold, 2005), and in the
sparsest constructions of networks.

Mark Kac asked in 1966 whether the frequencies of a drum, the eigenvalues of its Laplacian, determine its shape. For
graphs the answer is no (Exercise~\ref{spectral:ex:cospectral}), but the eigenvalues know a remarkable amount: the number of spanning
trees, the number of Eulerian circuits, how well the graph can be cut in two, how fast a random walk forgets where it started.
This chapter proves these facts. Tome~2, Chapter~8 studies the adjacency matrix itself (the friendship theorem, Moore graphs,
Hoffman's bound, Huang's theorem); here the main character is the \emph{Laplacian}.

\begin{think}
How many spanning trees does the complete graph $K_4$ have? Draw them all. Now guess how many spanning trees the cube graph
$Q_3$ has, with its $8$ vertices and $12$ edges.
\end{think}

\section{The Laplacian}

Let $G$ be a graph (multiple edges allowed, no loops) with vertex set $V=\{1,\dots,n\}$. Its \emph{adjacency matrix} $A$ has $a_{uv}$ equal to the
number of edges between $u$ and $v$, its \emph{degree matrix} $D$ is diagonal with the degrees $d(v)$, and its \emph{Laplacian} is
\[
L=D-A.
\]
Orient each edge arbitrarily and let $B$ be the $n\times m$ \emph{incidence matrix}: the column of an edge from $u$ to $v$ has $+1$ in row $u$, $-1$ in
row $v$, and $0$ elsewhere. Then $(BB^{\mathsf T})_{uv}=\sum_eb_{ue}b_{ve}$ is $d(u)$ if $u=v$ and minus the number of edges between $u$ and $v$ otherwise:
$L=BB^{\mathsf T}$, whatever the orientation. Hence, for $x\in\R^V$,
\[
x^{\mathsf T}Lx=\|B^{\mathsf T}x\|^2=\sum_{uv\in E}(x_u-x_v)^2.
\]
The Laplacian measures how much a function on the vertices varies along the edges: it is a discrete version of $\int|\nabla f|^2$, which
explains the name. By the spectral theorem (Tome~0, Theorem~7.7), $L$ has real eigenvalues
\[
0=\mu_1\le\mu_2\le\dots\le\mu_n
\]
and an orthonormal basis of eigenvectors $\varphi_1,\dots,\varphi_n$; they are non-negative since $x^{\mathsf T}Lx\ge0$, and $\mu_1=0$ since $L\mathbf1=0$, where $\mathbf1$ is the all-ones
vector. We can take $\varphi_1=\frac1{\sqrt n}\mathbf1$. For a $d$-regular graph, $L=dI-A$, and the eigenvalues of $L$ are $d-\lambda$ for the eigenvalues $\lambda$ of $A$.

\begin{lemma}[Rayleigh]\label{spectral:lem:rayleigh}
Let $M$ be a real symmetric matrix with eigenvalues $\nu_1\le\dots\le\nu_n$ and orthonormal eigenvectors $\psi_1,\dots,\psi_n$. If $x\ne0$ is orthogonal to
$\psi_1,\dots,\psi_{k-1}$, then $x^{\mathsf T}Mx\ge\nu_kx^{\mathsf T}x$, with equality for $x=\psi_k$. Similarly $x^{\mathsf T}Mx\le\nu_nx^{\mathsf T}x$ for every $x$. In particular
\[
\mu_2=\min\Bigl\{\frac{x^{\mathsf T}Lx}{x^{\mathsf T}x}:x\ne0,\ x\perp\mathbf1\Bigr\}.
\]
\end{lemma}
\begin{proof}
Write $x=\sum_ic_i\psi_i$; then $x^{\mathsf T}x=\sum_ic_i^2$ and $x^{\mathsf T}Mx=\sum_i\nu_ic_i^2$. If $c_1=\dots=c_{k-1}=0$, this is at least $\nu_k\sum_ic_i^2$; and it is always at most
$\nu_n\sum_ic_i^2$. For $L$, take $\psi_1=\varphi_1$, proportional to $\mathbf1$.
\end{proof}

Tome~2, Theorem~E.4 (Courant--Fischer) describes every eigenvalue in this way; we need only the second one.

\begin{proposition}\label{spectral:prop:components}
The multiplicity of the eigenvalue $0$ of $L$ is the number of connected components of $G$. In particular $\mu_2>0$ if and only if $G$ is connected.
\end{proposition}
\begin{proof}
$Lx=0$ implies $x^{\mathsf T}Lx=\sum_{uv\in E}(x_u-x_v)^2=0$, so $x$ is constant along every edge, hence on every component; conversely such $x$ satisfy $B^{\mathsf T}x=0$,
so $Lx=BB^{\mathsf T}x=0$. The kernel of $L$ is therefore the space of functions constant on each component, of dimension the number of components.
\end{proof}

Miroslav Fiedler called $\mu_2$ the \emph{algebraic connectivity} of $G$ (1973), and the eigenvectors of $\mu_2$ are now called \emph{Fiedler vectors}.
Section~\ref{spectral:sec:cheeger} shows that $\mu_2$ measures not only whether $G$ is connected but how well.

\begin{example}\label{spectral:ex:lapspectra}
From the adjacency spectra of Tome~2, Example~H.1, for regular graphs: the complete graph $K_n$ has Laplacian eigenvalues $0$ and $n$ (multiplicity $n-1$);
the cycle $C_n$ has $2-2\cos\frac{2\pi j}n$, $j=0,\dots,n-1$; the hypercube $Q_n$ has $2j$ with multiplicity $\binom nj$, $j=0,\dots,n$. The Petersen graph, strongly regular,
has adjacency eigenvalues $3$, $1$ (five times) and $-2$ (four times) by Tome~2, Theorem~H.4, hence Laplacian eigenvalues $0$, $2$ (five times), $5$ (four times).
\end{example}

\section{Counting spanning trees}

The cleanest route to Kirchhoff's theorem goes through directed graphs, where the theorem is due to Tutte (1948). Let $G$ be a directed graph on
$V=\{1,\dots,n\}$, with multiple arcs and loops allowed. Let $a_{uv}$ be the number of arcs from $u$ to $v$, $d^+(u)$ the out-degree of $u$ (loops
included), and $L=D^+-A$. Loops cancel on the diagonal, so they do not affect $L$. A \emph{spanning arborescence rooted at $w$} (or \emph{in-tree}) is a set of
$n-1$ arcs, one leaving each vertex $u\ne w$, such that following them from any vertex leads to $w$. For a matrix $M$ indexed by $V$, write $M_w$ for $M$ with the row
and column of $w$ deleted.

\begin{theorem}[matrix-tree theorem, directed form]\label{spectral:thm:dirmatrixtree}
The number $t_w(G)$ of spanning arborescences rooted at $w$ is $\det L_w$.
\end{theorem}
\begin{proof}
Row $u$ of $L$ is $\sum_{a}(e_u-e_{h(a)})$, the sum over the arcs $a$ leaving $u$, where $h(a)$ is the head of $a$ and $e_v$ the standard basis vector (a loop gives
$e_u-e_u=0$). Deleting the column of $w$ is linear, so the rows of $L_w$ are sums of the truncated vectors $\bar e_u-\bar e_{h(a)}$, where $\bar e_w=0$. The determinant is linear in each
row, so
\[
\det L_w=\sum_f\det M_f,
\]
the sum over all maps $f$ choosing an arc $f(u)$ leaving each vertex $u\ne w$, where $M_f$ has rows $\bar e_u-\bar e_{h(f(u))}$. The chosen arcs form a graph in which every $u\ne w$ has exactly one
outgoing arc. Two cases.

If following the chosen arcs from some vertex never reaches $w$, the path must eventually repeat a vertex: there is a cycle $u_1\to u_2\to\dots\to u_k\to u_1$ of chosen arcs, avoiding $w$ (a loop
is the case $k=1$). The sum of the rows $u_1,\dots,u_k$ of $M_f$ is $\sum_i(\bar e_{u_i}-\bar e_{u_{i+1}})=0$ (indices modulo $k$), so $\det M_f=0$.

Otherwise the chosen arcs form an arborescence rooted at $w$. Number the vertices $u\ne w$ by increasing distance to $w$ along the chosen arcs. The head of $f(u)$ is
$w$ or a vertex closer to $w$, hence earlier. So in this order, $M_f$ is lower triangular with $1$ on the diagonal, and $\det M_f=1$. The sum counts the arborescences.
\end{proof}

An undirected graph is turned into a directed one by replacing each edge $uv$ by the two arcs $u\to v$ and $v\to u$. The matrix $D^+-A$ is then the Laplacian of
the graph, and the arborescences rooted at $w$ are the spanning trees, each oriented towards $w$.

\begin{corollary}[Kirchhoff, 1847]\label{spectral:cor:kirchhoff}
For a graph $G$ on $n$ vertices and any vertex $w$, the number $\tau(G)$ of spanning trees equals $\det L_w$, and
\[
\tau(G)=\frac1n\mu_2\mu_3\cdots\mu_n.
\]
\end{corollary}
\begin{proof}
The first formula is Theorem~\ref{spectral:thm:dirmatrixtree} for the doubled graph. For the second, expand the characteristic polynomial: for any $n\times n$ matrix $M$, the coefficient of
$t$ in $\det(tI-M)$ is $(-1)^{n-1}\sum_w\det M_w$ (choose the entry $t$ in one diagonal position $w$, and the constant terms elsewhere). For $M=L$, $\det(tI-L)=\prod_i(t-\mu_i)=t\prod_{i\ge2}(t-\mu_i)$, whose
coefficient of $t$ is $(-1)^{n-1}\mu_2\cdots\mu_n$. Hence $\sum_w\det L_w=\mu_2\cdots\mu_n$, and each of the $n$ terms is $\tau(G)$.
\end{proof}

\begin{corollary}\label{spectral:cor:treecounts}
$\tau(K_n)=n^{n-2}$, and the hypercube has
\[
\tau(Q_n)=2^{2^n-n-1}\prod_{j=1}^nj^{\binom nj}.
\]
\end{corollary}
\begin{proof}
By Example~\ref{spectral:ex:lapspectra}, $\tau(K_n)=\frac1nn^{n-1}$, which is Cayley's formula again (Tome~3, Chapter~12), in Borchardt's way of 1860. For the hypercube,
$\tau(Q_n)=2^{-n}\prod_{j=1}^n(2j)^{\binom nj}$, and $\sum_{j=1}^n\binom nj=2^n-1$.
\end{proof}

So $Q_3$ has $2^4\cdot1\cdot2^3\cdot3=384$ spanning trees, and the Petersen graph $\frac1{10}2^5\cdot5^4=2000$. Proofs of the hypercube formula without eigenvalues are much harder: Olivier Bernardi found a combinatorial one only in 2012, and a simple bijection is still sought.
It is a case where eigenvalues see further than combinatorics.

\begin{figure}[ht]
\centering
\fig{spectral-ch09-trees}
\caption{The $16=4^2$ spanning trees of $K_4$. Twelve are paths and four are stars; Kirchhoff's determinant counts them all at once.}
\label{spectral:fig:k4trees}
\end{figure}

\section{Eulerian circuits}

The directed theorem has a beautiful partner. A directed graph is \emph{Eulerian} if $d^+(v)=d^-(v)$ for every vertex and every vertex can be reached from every
other (Hierholzer's argument splices closed tours until every arc is used, giving an Eulerian circuit). How many Eulerian circuits are there? We count them as
cyclic sequences of arcs, so that rotating a circuit does not give a new one.

\begin{theorem}[BEST theorem: de Bruijn, van Aardenne-Ehrenfest, Smith, Tutte, 1941--1951]\label{spectral:thm:best}
Let $G$ be an Eulerian directed graph and $w$ a vertex. The number of Eulerian circuits of $G$ is
\[
\mathrm{ec}(G)=t_w(G)\prod_{v\in V}\bigl(d^+(v)-1\bigr)!.
\]
\end{theorem}
\begin{proof}
Fix an arc $e$ leaving $w$. Each Eulerian circuit can be rotated in exactly one way into an \emph{Eulerian trail} starting with $e$, a sequence of all arcs starting with $e$, each arc
starting where the previous one ends, and ending at $w$. We count these trails. For such a trail $W$ and $v\ne w$, let $\ell(v)$ be the arc by which $W$ leaves $v$ for the last time.

\emph{The last exits form an arborescence rooted at $w$.} Each $v\ne w$ has exactly one arc $\ell(v)$. If $\ell(v)$ goes from $v$ to $u\ne w$, then after traversing $\ell(v)$ the trail is at $u$, and it
must leave $u$ later, since it ends at $w\ne u$: the last exit from $u$ comes after the last exit from $v$. Along the arcs $\ell$, the times of last exits increase, so there is no cycle, and
following the arcs $\ell$ from any vertex leads to $w$.

\emph{The trail is determined by the arborescence and some orders.} Record, for each $v\ne w$, the order in which $W$ uses the arcs leaving $v$ other than $\ell(v)$, which are used before
$\ell(v)$; and for $w$, the order in which it uses the arcs leaving $w$ other than $e$, which comes first. The number of choices of these orders is $\prod_v(d^+(v)-1)!$. Conversely, given an
arborescence $T$ rooted at $w$ and such orders, the trail is rebuilt by the following rule: start with $e$; each time you are at a vertex, leave by its next unused arc in the prescribed
order, where for $v\ne w$ the arc of $T$ comes last. This rule reproduces $W$ from its own data. It remains to show that for any $T$ and any orders, the rule produces an Eulerian trail
(whose last exits are then the arcs of $T$, and whose orders are the given ones).

The walk can only stop at a vertex with no unused arc leaving it. When it is at $v\ne w$, it has entered $v$ once more than it has left it, so it has used at most $d^-(v)-1=d^+(v)-1$ arcs leaving $v$:
it never stops at $v\ne w$. So it stops at $w$, having used all the arcs leaving $w$, and since it started and ended at $w$, it has entered $w$ as often as it has left it: all $d^-(w)=d^+(w)$ arcs
entering $w$ are used. Suppose some arc leaving a vertex $v$ is unused. Then the arc of $T$ at $v$, which comes last, is unused; it goes to a vertex $u$. If $u\ne w$, then $u$ has an unused entering arc,
so it was entered fewer than $d^-(u)$ times and left as often as entered, fewer than $d^+(u)$ times: the arc of $T$ at $u$ is unused too. Following $T$ we reach $w$ through an unused arc entering $w$,
a contradiction. So every arc is used: the walk is an Eulerian trail.
\end{proof}

Since $\mathrm{ec}(G)$ and the product do not depend on $w$, neither does $t_w(G)$ for an Eulerian graph, a fact not obvious at all from the definition. Here is the application that de Bruijn
had in mind. A binary \emph{de Bruijn sequence of order $n$} is a cyclic word of length $2^n$ in which every binary word of length $n$ appears exactly once as a window
(the Eulerian circuit construction below proves existence). How many are there?

\begin{corollary}[Flye Sainte-Marie, 1894; de Bruijn, 1946]\label{spectral:cor:debruijncount}
There are exactly $2^{2^{n-1}-n}$ binary de Bruijn sequences of order $n$, counted up to rotation.
\end{corollary}
\begin{proof}
For $n=1$ there is one, $01$. For $n\ge2$, the usual overlapping-word construction matches de Bruijn sequences up to rotation with the Eulerian circuits of the de Bruijn graph $D(2,n)$, whose $N=2^{n-1}$
vertices are the words of length $n-1$, with an arc from $a_1\cdots a_{n-1}$ to $a_2\cdots a_{n-1}b$ for each $b\in\{0,1\}$ (two of the arcs are loops). Every vertex has out-degree $2$, so by
Theorem~\ref{spectral:thm:best} the number of circuits is $t_w\cdot1!^N=t_w$, and it remains to count arborescences.

There is exactly one walk of length $n-1$ from any vertex to any other (shift in the letters of the target), so $A^{n-1}=J$, the all-ones matrix. The eigenvalues of $J$ are $N$ once and $0$
with multiplicity $N-1$, and they are the $(n-1)$-th powers of the eigenvalues of $A$, with multiplicities (triangularise $A$ over $\C$). Since $A\mathbf1=2\mathbf1$ and $2^{n-1}=N$, the eigenvalues
of $A$ are $2$ once and $0$ with multiplicity $N-1$, and those of $L=2I-A$ are $0$ once and $2$ with multiplicity $N-1$. As in Corollary~\ref{spectral:cor:kirchhoff}, $\sum_w\det L_w=2^{N-1}$; all terms are equal
to $t_w$, so $t_w=\frac{2^{N-1}}N=2^{2^{n-1}-1-(n-1)}$.
\end{proof}

For $n=3$ there are $2^{4-3}=2$ of them, $00010111$ and $00011101$; for $n=4$ there are $16$, and for $n=6$ already $2^{26}$. Figure~\ref{spectral:fig:best} shows the bijection at work for $n=3$.

\begin{figure}[ht]
\centering
\fig{spectral-ch09-best}
\caption{The de Bruijn graph $D(2,3)$. The thick arcs form an arborescence rooted at $00$ (the last exits). Starting with the loop at $00$ and leaving each vertex by its
non-tree arc first, the rule of the BEST theorem walks $00\to00\to01\to11\to11\to10\to01\to10\to00$, and the first letters of the arcs spell the de Bruijn sequence $00011101$.}
\label{spectral:fig:best}
\end{figure}

\section{Cutting a graph: Cheeger's inequality}\label{spectral:sec:cheeger}

How well connected is a connected graph? A natural measure is the cheapest way to cut off a piece of it, relative to the size of the piece. Let $G$ be $d$-regular with $n$
vertices. For $S\subseteq V$, write $\bar S=V\setminus S$ and $e(S,\bar S)$ for the number of edges between $S$ and $\bar S$. The \emph{edge expansion}, or \emph{Cheeger constant}, of $G$ is
\[
h(G)=\min\Bigl\{\frac{e(S,\bar S)}{|S|}:S\ne\varnothing,\ |S|\le\frac n2\Bigr\}.
\]
Computing $h(G)$ exactly is NP-hard in general, a computational hardness property. The second eigenvalue gives it up to a square root, and its eigenvector gives a good cut.

\begin{theorem}[Cheeger's inequality: Dodziuk, 1984; Alon and Milman, 1985]\label{spectral:thm:cheeger}
For every $d$-regular graph,
\[
\frac{\mu_2}2\le h(G)\le\sqrt{2d\mu_2}.
\]
\end{theorem}
\begin{proof}
\emph{Lower bound.} Let $S$ achieve $h(G)$, $s=|S|\le\frac n2$, and $x=\mathbf1_S-\frac sn\mathbf1$, which is orthogonal to $\mathbf1$. An edge contributes $(x_u-x_v)^2=1$ if it crosses between
$S$ and $\bar S$ and $0$ otherwise, so $x^{\mathsf T}Lx=e(S,\bar S)$, while $x^{\mathsf T}x=s(1-\frac sn)^2+(n-s)\frac{s^2}{n^2}=\frac{s(n-s)}n$. By Lemma~\ref{spectral:lem:rayleigh},
\[
\mu_2\le\frac{n\,e(S,\bar S)}{s(n-s)}\le\frac{2e(S,\bar S)}s=2h(G),
\]
as $n-s\ge\frac n2$.

\emph{Upper bound.} Let $x$ be an eigenvector of $\mu_2$. It is orthogonal to $\mathbf1$ and non-zero, so it has positive and negative entries; replacing $x$ by $-x$ if necessary, the set
$V^+=\{v:x_v>0\}$ has at most $\frac n2$ elements. Let $y_v=\max(x_v,0)$.

First, $y^{\mathsf T}Ly\le\mu_2\,y^{\mathsf T}y$. For $v\in V^+$, $(Ly)_v=dy_v-\sum_{u\sim v}y_u\le dx_v-\sum_{u\sim v}x_u=(Lx)_v=\mu_2x_v=\mu_2y_v$, since $y_u\ge x_u$ and $y_v=x_v$. So
$y^{\mathsf T}Ly=\sum_{v\in V^+}y_v(Ly)_v\le\mu_2\sum_vy_v^2$.

Second, for $t\ge0$ let $S_t=\{v:y_v^2>t\}\subseteq V^+$, so $|S_t|\le\frac n2$ and $e(S_t,\bar S_t)\ge h(G)|S_t|$. An edge $uv$ crosses the cut $(S_t,\bar S_t)$ exactly for $t$ between $y_u^2$ and $y_v^2$,
an interval of length $|y_u^2-y_v^2|$. Integrating over $t$,
\[
\sum_{uv\in E}|y_u^2-y_v^2|=\int_0^\infty e(S_t,\bar S_t)\dd t\ge h(G)\int_0^\infty|S_t|\dd t=h(G)\sum_vy_v^2,
\]
since a vertex $v$ lies in $S_t$ for $t$ in an interval of length $y_v^2$.

Third, by the Cauchy--Schwarz inequality and $(a+b)^2\le2a^2+2b^2$,
\[
\begin{aligned}\sum_{uv\in E}|y_u^2-y_v^2|&=\sum_{uv\in E}|y_u-y_v|(y_u+y_v)\eqbr\le\sqrt{y^{\mathsf T}Ly}\,\sqrt{\textstyle\sum_{uv\in E}2(y_u^2+y_v^2)}=\sqrt{y^{\mathsf T}Ly}\,\sqrt{2d\textstyle\sum_vy_v^2},\end{aligned}
\]
as each vertex is in $d$ edges. Combining the three steps, $h(G)\sum_vy_v^2\le\sqrt{2d\mu_2}\sum_vy_v^2$, and $y\neq0$.
\end{proof}

The proof is an algorithm, called \emph{spectral partitioning}. The integral inequality $\int_0^\infty e(S_t,\bar S_t)\dd t\le\sqrt{2d\mu_2}\int_0^\infty|S_t|\dd t$ shows that some $t$ satisfies
$e(S_t,\bar S_t)\le\sqrt{2d\mu_2}\,|S_t|$, and the sets $S_t$ are the $n$ or fewer \emph{sweep sets}: sort the vertices by the value of the Fiedler vector and try every threshold. In 2000 Jianbo Shi and
Jitendra Malik used a normalised version of this idea to segment images, the pixels being vertices joined according to their similarity, and \emph{spectral clustering} became a standard tool of
data analysis.

Both inequalities are tight up to constants. For the hypercube $Q_n$, $\mu_2=2$ and $h(Q_n)=1$ (cut along a coordinate; that nothing is better is the edge-isoperimetric inequality, a classical
exercise of combinatorics), so the left side is attained. For the cycle, $\mu_2\approx\frac{4\pi^2}{n^2}$ while $h(C_n)\approx\frac4n$, so the square root on the right is necessary (Exercise~\ref{spectral:ex:cheegercycle}).

\begin{figure}[ht]
\centering
\fig{spectral-ch09-fiedler}
\caption{Top: a graph made of two dense clusters of $11$ vertices joined by three edges (red). Bottom: the values of its Fiedler vector $\varphi_2$, one dot per vertex. The best
sweep set (dashed threshold) is exactly one cluster (blue): the cut crosses only the three bridges.}
\label{spectral:fig:fiedler}
\end{figure}

\section{Expanders}

For a $d$-regular graph, let $d=\lambda_1\ge\lambda_2\ge\dots\ge\lambda_n$ be the eigenvalues of $A$ (so $\mu_i=d-\lambda_i$), and $\lambda=\max(|\lambda_2|,|\lambda_n|)$, the largest absolute value of an
eigenvalue other than $d$. A $d$-regular graph on $n$ vertices with this $\lambda$ is called an \emph{$(n,d,\lambda)$-graph}. Cheeger's inequality says that $h(G)\ge\frac{d-\lambda_2}2$: if $\lambda$ is small compared
with $d$, every set of at most half the vertices has many edges leaving it. A family of $d$-regular graphs with $n\to\infty$ and $\lambda\le(1-\varepsilon)d$ is a family of \emph{expanders}. The
next lemma says that in such a graph, edges are spread as if at random.

\begin{lemma}[expander mixing lemma, Alon and Chung, 1988]\label{spectral:lem:mixing}
Let $G$ be an $(n,d,\lambda)$-graph and $S,T\subseteq V$. Let $e(S,T)$ be the number of pairs $(s,t)\in S\times T$ with $s\sim t$. Then
\[
\Bigl|e(S,T)-\frac{d|S||T|}n\Bigr|\le\lambda\sqrt{|S||T|\Bigl(1-\frac{|S|}n\Bigr)\Bigl(1-\frac{|T|}n\Bigr)}\le\lambda\sqrt{|S||T|}.
\]
\end{lemma}
\begin{proof}
Write $\mathbf1_S=\frac{|S|}n\mathbf1+x$ and $\mathbf1_T=\frac{|T|}n\mathbf1+y$, with $x,y\perp\mathbf1$. Then $e(S,T)=\mathbf1_S^{\mathsf T}A\mathbf1_T$, and since $A\mathbf1=d\mathbf1$ and $A$ maps $\mathbf1^\perp$ to itself,
\[
e(S,T)=\frac{|S||T|}{n^2}\,\mathbf1^{\mathsf T}A\mathbf1+x^{\mathsf T}Ay=\frac{d|S||T|}n+x^{\mathsf T}Ay.
\]
Expand $x=\sum_{i\ge2}c_i\psi_i$ and $y=\sum_{i\ge2}c'_i\psi_i$ in an orthonormal eigenbasis of $A$ with $\psi_1\propto\mathbf1$: $|x^{\mathsf T}Ay|=|\sum_{i\ge2}\lambda_ic_ic'_i|\le\lambda\|x\|\|y\|$ by the Cauchy--Schwarz
inequality. Finally $\|x\|^2=\|\mathbf1_S\|^2-\frac{|S|^2}n=|S|(1-\frac{|S|}n)$, and likewise for $y$.
\end{proof}

For instance an independent set $S$ has $e(S,S)=0$, so $\frac{d|S|^2}n\le\lambda|S|(1-\frac{|S|}n)$, that is $|S|\le\frac{n\lambda}{d+\lambda}$ (compare Hoffman's bound, Tome~2, Theorem~H.6). The lemma also
bounds the diameter (Exercise~\ref{spectral:ex:diameter}) and shows that random walks forget their starting point exponentially fast (Exercise~\ref{spectral:ex:walkmixing}). Good expanders want $\lambda$ as small as
possible. How small can it be?

\begin{theorem}[Alon and Boppana, 1986; Nilli's proof, 1991]\label{spectral:thm:alonboppana}
Let $G$ be $d$-regular, $d\ge2$, with two edges $\{v_1,v_2\}$ and $\{u_1,u_2\}$ whose endpoints are at distance at least $2k+2$ from each other. Then
\[
\lambda_2\ge2\sqrt{d-1}-\frac{2\sqrt{d-1}-1}{k+1}.
\]
\end{theorem}
\begin{proof}
We build a test vector for Lemma~\ref{spectral:lem:rayleigh} that looks, near each of the two edges, like the eigenvector of an infinite $d$-regular tree. Let $V_i$ be the set of vertices at distance
exactly $i$ from $\{v_1,v_2\}$, for $0\le i\le k$, and $a_i=(d-1)^{-i/2}$. Define $f_x=a_i$ for $x\in V_i$, $i\le k$, and $f_x=0$ elsewhere.

Every vertex of $V_i$ with $i\ge1$ has a neighbour in $V_{i-1}$, and each of $v_1,v_2$ has the other as a neighbour in $V_0$; so each vertex of $V_i$ has at most $d-1$ neighbours in $V_{i+1}$, and
$|V_{i+1}|\le(d-1)|V_i|$. Edges inside some $V_i$, or outside $V_0\cup\dots\cup V_{k+1}$, contribute nothing to $f^{\mathsf T}Lf=\sum_{xy\in E}(f_x-f_y)^2$. Edges between $V_i$ and $V_{i+1}$ number at most
$(d-1)|V_i|$, so
\[
f^{\mathsf T}Lf\le\sum_{i=0}^{k-1}(d-1)|V_i|(a_i-a_{i+1})^2+(d-1)|V_k|a_k^2.
\]
Now $(d-1)(a_i-a_{i+1})^2=a_i^2(\sqrt{d-1}-1)^2=a_i^2(d-2\sqrt{d-1})$, and $(d-1)a_k^2=a_k^2(d-2\sqrt{d-1})+a_k^2(2\sqrt{d-1}-1)$. Hence
\[
f^{\mathsf T}Lf\le(d-2\sqrt{d-1})\sum_{i=0}^k|V_i|a_i^2+(2\sqrt{d-1}-1)|V_k|a_k^2.
\]
The numbers $|V_i|a_i^2$ do not increase with $i$, since $|V_{i+1}|a_{i+1}^2\le(d-1)|V_i|\frac{a_i^2}{d-1}$; so the last one is at most their average, $|V_k|a_k^2\le\frac1{k+1}\sum_i|V_i|a_i^2=\frac{f^{\mathsf T}f}{k+1}$. With
$\theta=d-2\sqrt{d-1}+\frac{2\sqrt{d-1}-1}{k+1}$ we get $f^{\mathsf T}Lf\le\theta f^{\mathsf T}f$.

Define $g$ in the same way around $\{u_1,u_2\}$; also $g^{\mathsf T}Lg\le\theta g^{\mathsf T}g$. The supports of $f$ and $g$ are at distance at least $2k+2-2k=2$, so they are disjoint and no edge joins them.
Choose $c$ with $h=f-cg\perp\mathbf1$ (possible as $\sum g>0$). Then $h\ne0$, and since no edge meets both supports, $h^{\mathsf T}Lh=f^{\mathsf T}Lf+c^2g^{\mathsf T}Lg\le\theta(f^{\mathsf T}f+c^2g^{\mathsf T}g)=\theta h^{\mathsf T}h$. By
Lemma~\ref{spectral:lem:rayleigh}, $\mu_2\le\theta$, that is $\lambda_2=d-\mu_2\ge d-\theta$.
\end{proof}

A $d$-regular graph with $n$ vertices has at most $2(1+(d-1)+\dots+(d-1)^r)$ vertices within distance $r$ of a given edge, so when $n\to\infty$ with $d$ fixed, there are edges at distance $2k+2$ for larger
and larger $k$: $\lambda_2\ge2\sqrt{d-1}-o(1)$. The number $2\sqrt{d-1}$ is the spectral radius of the infinite $d$-regular tree, which every large $d$-regular graph resembles locally. Graphs with
$\lambda\le2\sqrt{d-1}$ are called \emph{Ramanujan graphs}: Lubotzky, Phillips and Sarnak (1988), and independently Margulis, built them for $d-1$ prime using Ramanujan's conjecture on the
coefficients of modular forms (proved by Deligne), and Marcus, Spielman and Srivastava (2013) proved that bipartite Ramanujan graphs exist for every degree, with a new method of ``interlacing
families'' of polynomials. Tome~11 tells this story. And Joel Friedman proved in 2008 that a \emph{random} $d$-regular graph has $\lambda\le2\sqrt{d-1}+\varepsilon$ with probability tending to $1$: almost all
regular graphs are almost optimal expanders, while building one explicitly is hard.

\begin{beauty}
Three numbers attached to a matrix, the product of the non-zero eigenvalues, the second smallest eigenvalue, and the gap to the spectral radius of a tree, answer three questions
that seem to have nothing to do with linear algebra: how many ways a network can be kept connected with the fewest wires, how cheaply it can be cut, and how fast information
spreads in it. The matrix $L=BB^{\mathsf T}$, a square of the incidence matrix, is all it takes. Kac's drum could not be heard completely, but it can be heard surprisingly well.
\end{beauty}

\section{Exercises}

\begin{exercise}[$\star$]\label{spectral:ex:kmntrees}
Find the Laplacian eigenvalues of the complete bipartite graph $K_{m,n}$, and deduce that $\tau(K_{m,n})=m^{n-1}n^{m-1}$.
\end{exercise}

\begin{exercise}[$\star$]\label{spectral:ex:delcontrees}
Show that $\tau(G)=\tau(G-e)+\tau(G/e)$ for every edge $e$ that is not a loop, where $G/e$ (contraction) may have multiple edges. Check $\tau(K_4)=16$ in this way.
\end{exercise}

\begin{exercise}[$\star$]\label{spectral:ex:cospectral}
Show that the star $K_{1,4}$ and the disjoint union of a square $C_4$ and an isolated vertex have the same adjacency eigenvalues, but different Laplacian eigenvalues. Which of the two
spectra knows whether a graph is connected?
\end{exercise}

\begin{exercise}[$\star\star$]\label{spectral:ex:debruijnq}
Show that the number of de Bruijn sequences of order $n$ over an alphabet of $q$ letters, up to rotation, is $\frac{(q!)^{q^{n-1}}}{q^n}$.
\end{exercise}

\begin{exercise}[$\star$]\label{spectral:ex:arborescences}
Let $G$ be the directed graph on $\{1,2,3\}$ with arcs $1\to2$, $2\to3$, $3\to1$, $1\to3$, $3\to2$, $2\to1$, $1\to2$ (so two arcs from $1$ to $2$). Compute $t_w(G)$ for each $w$ with
Theorem~\ref{spectral:thm:dirmatrixtree}, and check by listing the arborescences. Is $G$ Eulerian?
\end{exercise}

\begin{exercise}[$\star$]\label{spectral:ex:cheegercycle}
Show that $h(C_n)=\frac2{\lfloor n/2\rfloor}$ and $\mu_2(C_n)=2-2\cos\frac{2\pi}n$. Conclude that in Theorem~\ref{spectral:thm:cheeger} the upper bound cannot be replaced by $c\,\mu_2$ for any constant $c$.
\end{exercise}

\begin{exercise}[$\star\star$]\label{spectral:ex:diameter}
Let $G$ be a connected $(n,d,\lambda)$-graph with $\lambda<d$. Show that $(A^t)_{uv}\ge\frac{d^t}n-\lambda^t$ for all vertices $u,v$, and deduce that the diameter of $G$ is at most
$\bigl\lfloor\frac{\ln n}{\ln(d/\lambda)}\bigr\rfloor+1$.
\end{exercise}

\begin{exercise}[$\star$]\label{spectral:ex:walkmixing}
A random walk on a $d$-regular graph moves from its current vertex to a uniform random neighbour. Let $p_t$ be the distribution after $t$ steps and $u=\frac1n\mathbf1$. Show that $p_{t+1}=\frac1dAp_t$
and $\|p_t-u\|\le(\frac\lambda d)^t$ for an $(n,d,\lambda)$-graph.
\end{exercise}

\begin{exercise}[$\star$]\label{spectral:ex:petersenalpha}
Use Lemma~\ref{spectral:lem:mixing} to show that the Petersen graph has no independent set of $5$ vertices. Find one of $4$.
\end{exercise}

\begin{exercise}[$\star\star$]\label{spectral:ex:fiedlerpath}
Show that the eigenvalues of the Laplacian of the path $P_n$ are $2-2\cos\frac{j\pi}n$, $j=0,\dots,n-1$, with eigenvectors $x_v=\cos\frac{j\pi(2v-1)}{2n}$, $v=1,\dots,n$. What cut does the sweep of
the Fiedler vector find?
\end{exercise}

\begin{exercise}[$\star$]\label{spectral:ex:moorebound}
(a) Show that at most $2\bigl(1+(d-1)+\dots+(d-1)^r\bigr)$ vertices of a $d$-regular graph are within distance $r$ of a given edge, with equality when the graph has girth at least $2r+2$.
(b) Deduce that a $d$-regular graph with more than $2\bigl(1+(d-1)+\dots+(d-1)^{2k+2}\bigr)$ vertices has two edges at distance at least $2k+2$, so that Theorem~\ref{spectral:thm:alonboppana} applies.
\end{exercise}

\section{Hints}

\begin{hint}{spectral:ex:kmntrees}
Look for eigenvectors constant on each side, and eigenvectors supported on one side with sum $0$.
\end{hint}
\begin{hint}{spectral:ex:delcontrees}
Split the spanning trees according to whether they contain $e$.
\end{hint}
\begin{hint}{spectral:ex:cospectral}
Both have $\operatorname{tr}A^2=8$; compute the characteristic polynomials.
\end{hint}
\begin{hint}{spectral:ex:debruijnq}
In $D(q,n)$, $A^{n-1}=J$ again.
\end{hint}
\begin{hint}{spectral:ex:arborescences}
$L=D^+-A$ is $3\times3$; delete a row and a column.
\end{hint}
\begin{hint}{spectral:ex:cheegercycle}
A set of $s$ vertices of a cycle has at least $2$ boundary edges, exactly $2$ for an arc.
\end{hint}
\begin{hint}{spectral:ex:diameter}
Use the spectral decomposition $A^t=\sum_i\lambda_i^t\psi_i\psi_i^{\mathsf T}$ and $\sum_i\psi_i(u)^2=1$.
\end{hint}
\begin{hint}{spectral:ex:walkmixing}
$p_t-u$ is orthogonal to $\mathbf1$.
\end{hint}
\begin{hint}{spectral:ex:petersenalpha}
$\lambda=2$ for the Petersen graph.
\end{hint}
\begin{hint}{spectral:ex:fiedlerpath}
Check $(Lx)_v=\mu x_v$ at the inner vertices with $\cos(a-b)+\cos(a+b)=2\cos a\cos b$, then at the ends.
\end{hint}
\begin{hint}{spectral:ex:moorebound}
Count the vertices at distance exactly $i$, as in the proof of Theorem~\ref{spectral:thm:alonboppana}.
\end{hint}

\section{Solutions}

\begin{solution}{spectral:ex:kmntrees}
\idea Symmetry suggests vectors constant on each side, and vectors that the matrix $J$ kills.
\writeup Let the sides be $X$ ($|X|=m$, degree $n$) and $Y$ ($|Y|=n$, degree $m$). A vector supported on $X$ with sum $0$ satisfies $Lx=nx$: on $X$, $(Lx)_v=nx_v$, and on $Y$, $(Lx)_y=-\sum_{v\in X}x_v=0$. This gives
the eigenvalue $n$ with multiplicity $m-1$, and likewise $m$ with multiplicity $n-1$. Vectors constant on each side, $x=\alpha\mathbf1_X+\beta\mathbf1_Y$, give $Lx=n(\alpha-\beta)\mathbf1_X+m(\beta-\alpha)\mathbf1_Y$: eigenvalues
$0$ (for $\alpha=\beta$) and $m+n$ (for $\alpha=n$, $\beta=-m$). That is $m+n$ eigenvalues in all. By Corollary~\ref{spectral:cor:kirchhoff}, $\tau=\frac1{m+n}n^{m-1}m^{n-1}(m+n)=m^{n-1}n^{m-1}$.
\end{solution}

\begin{solution}{spectral:ex:delcontrees}
\idea A tree either avoids $e$, and is a tree of $G-e$, or contains $e$, and becomes a tree of $G/e$ after contraction.
\writeup Spanning trees of $G$ without $e$ are exactly the spanning trees of $G-e$. Contracting $e=uv$ in a spanning tree containing $e$ gives a spanning tree of $G/e$ (the edges of $G/e$ are those of $G$
other than $e$, with $u$ and $v$ identified; parallel copies of $e$ become loops, which no tree uses), and conversely adding $e$ back to a spanning tree of $G/e$ gives a spanning tree of $G$ containing $e$.
For $K_4$: $K_4-e$ is a square with one diagonal, whose $8$ spanning trees are the $4$ paths around the square plus the $4$ trees using the diagonal and two of the four sides that do not form a
triangle with it ($\binom42$ choices minus $2$ triangles). $K_4/e$ has three vertices $w,a,b$ with double edges $wa$ and $wb$ and a single edge $ab$: its trees are pairs of edges not both between the same
two vertices, $2\cdot2+2\cdot1+2\cdot1=8$. Total $16$.
\end{solution}

\begin{solution}{spectral:ex:cospectral}
\idea Two small graphs with the same numbers of vertices, edges, triangles and closed walks of every length.
\writeup The characteristic polynomial of the adjacency matrix of the star is $t^3(t^2-4)$ (its rank is $2$ and $\operatorname{tr}A^2=8$, $\operatorname{tr}A=0$): eigenvalues $\pm2$ and $0,0,0$. The square has eigenvalues $2,0,0,-2$ (Tome~2, Example~H.1), and the isolated vertex
adds $0$. So both have spectrum $\{2,0,0,0,-2\}$. The Laplacians differ: the star has $0,1,1,1,5$ (it is $K_{1,4}$, Exercise~\ref{spectral:ex:kmntrees}), the other graph $0,0,2,2,4$. By Proposition~\ref{spectral:prop:components} the
Laplacian spectrum sees that the second graph is disconnected; the adjacency spectrum cannot.
\end{solution}

\begin{solution}{spectral:ex:debruijnq}
\idea The same computation as Corollary~\ref{spectral:cor:debruijncount}, with factorials that no longer equal $1$.
\writeup $D(q,n)$ has $N=q^{n-1}$ vertices, each of out-degree $q$, and exactly one walk of length $n-1$ between any two vertices, so $A^{n-1}=J$. As before the eigenvalues of $A$ are $q$ once and $0$ otherwise
(as $q^{n-1}=N$), and those of $L=qI-A$ are $0$ once and $q$ with multiplicity $N-1$. So $Nt_w=q^{N-1}$ and $t_w=q^{N-n}$. By Theorem~\ref{spectral:thm:best}, $\mathrm{ec}=q^{N-n}((q-1)!)^N=\frac{(q!)^N}{q^n}$.
\end{solution}

\begin{solution}{spectral:ex:arborescences}
\idea A $2\times2$ determinant for each root, then a list.
\writeup The out-degrees are $d^+(1)=3$ (arcs to $2,3,2$), $d^+(2)=2$, $d^+(3)=2$, and
\[
L=\begin{pmatrix}3&-2&-1\\-1&2&-1\\-1&-1&2\end{pmatrix}.
\]
Then $t_1=\det\begin{pmatrix}2&-1\\-1&2\end{pmatrix}=3$, $t_2=\det\begin{pmatrix}3&-1\\-1&2\end{pmatrix}=5$, $t_3=\det\begin{pmatrix}3&-2\\-1&2\end{pmatrix}=4$. For instance, rooted at $1$: choose the arcs out of $2$ and $3$ so that no cycle $2\to3\to2$
appears: $(2\to1,3\to1)$, $(2\to1,3\to2)$, $(2\to3,3\to1)$, three arborescences. The in-degree of $2$ is $3$ while its out-degree is $2$: $G$ is not Eulerian, and indeed $t_w$ depends on $w$.
\end{solution}

\begin{solution}{spectral:ex:cheegercycle}
\idea Cutting an arc of the cycle costs two edges; the eigenvalue is tiny because a slowly varying cosine has small energy.
\writeup A non-empty proper subset $S$ of the cycle is a union of $r\ge1$ arcs, with $2r\ge2$ boundary edges, so $\frac{e(S,\bar S)}{|S|}\ge\frac2{|S|}\ge\frac2{\lfloor n/2\rfloor}$, with equality for an arc of $\lfloor\frac n2\rfloor$ vertices. The
eigenvalues come from Example~\ref{spectral:ex:lapspectra}, the smallest non-zero one for $j=1$. As $n\to\infty$, $\mu_2\sim\frac{4\pi^2}{n^2}$ while $h(C_n)\sim\frac4n$: the ratio $\frac h{\mu_2}\sim\frac n{\pi^2}$ is unbounded, while
$\frac h{\sqrt{2d\mu_2}}\to\frac1\pi$ ($d=2$). So no bound $h\le c\mu_2$ holds, and the square root is the right order.
\end{solution}

\begin{solution}{spectral:ex:diameter}
\idea Count walks: the main term $\frac{d^t}n$ from the eigenvalue $d$ beats the error from the others once $t$ is large, and then every pair of vertices is joined by a walk of length $t$.
\writeup With an orthonormal eigenbasis, $\psi_1=\frac1{\sqrt n}\mathbf1$, $(A^t)_{uv}=\sum_i\lambda_i^t\psi_i(u)\psi_i(v)=\frac{d^t}n+\sum_{i\ge2}\lambda_i^t\psi_i(u)\psi_i(v)$. The rows of the orthogonal matrix $(\psi_i(u))$ have norm $1$, so
by Cauchy--Schwarz $\sum_{i\ge2}|\psi_i(u)\psi_i(v)|\le1$, and $(A^t)_{uv}\ge\frac{d^t}n-\lambda^t$. If $(\frac d\lambda)^t>n$, this is positive, so there is a walk of length $t$ from $u$ to $v$, and the distance is
at most $t$. Note that $\lambda>0$, as the eigenvalues add up to $\operatorname{tr}A=0$. The integer $t=\lfloor\frac{\ln n}{\ln(d/\lambda)}\rfloor+1$ exceeds $\frac{\ln n}{\ln(d/\lambda)}$, so $(\frac d\lambda)^t>n$ and the diameter is at most $t$. Expanders have logarithmic diameter.
\end{solution}

\begin{solution}{spectral:ex:walkmixing}
\idea The uniform distribution is the eigenvector of the top eigenvalue, and everything orthogonal to it shrinks by a factor $\frac\lambda d$ per step.
\writeup $p_{t+1}(v)=\sum_{u\sim v}\frac1dp_t(u)$, that is $p_{t+1}=\frac1dAp_t$ ($A$ is symmetric). Also $\frac1dAu=u$. So $p_{t+1}-u=\frac1dA(p_t-u)$, and $p_t-u$ has total mass $0$, so it lies in $\mathbf1^\perp$, where
$\frac1dA$ has eigenvalues of absolute value at most $\frac\lambda d$: $\|p_{t+1}-u\|\le\frac\lambda d\|p_t-u\|$. By induction $\|p_t-u\|\le(\frac\lambda d)^t\|p_0-u\|$, and $\|p_0-u\|^2=1-\frac1n\le1$ for a walk started at a vertex. A bipartite
graph has $\lambda_n=-d$, and the walk indeed never mixes: it alternates between the sides.
\end{solution}

\begin{solution}{spectral:ex:petersenalpha}
\idea The mixing lemma with $S=T$ independent.
\writeup The Petersen graph is a $(10,3,2)$-graph (eigenvalues $3,1,-2$). An independent set has $|S|\le\frac{n\lambda}{d+\lambda}=\frac{20}5=4$. In the usual drawing as the $2$-subsets of $\{1,\dots,5\}$, adjacent when disjoint,
the four $2$-sets containing $1$ pairwise intersect, so they are independent. Here the mixing bound is as good as Hoffman's.
\end{solution}

\begin{solution}{spectral:ex:fiedlerpath}
\idea Cosines satisfy the three-term recurrence, and the half-integer shift makes the ends behave like reflecting walls.
\writeup Let $\theta=\frac{j\pi}n$ and $x_v=\cos((v-\frac12)\theta)$. For $1<v<n$, $(Lx)_v=2x_v-x_{v-1}-x_{v+1}=2x_v-2\cos\theta\,x_v=(2-2\cos\theta)x_v$. At $v=1$, $(Lx)_1=x_1-x_2$, and the formula for the
inner vertices would give $2x_1-x_0-x_2$; they agree because $x_0=\cos(-\frac\theta2)=x_1$. Similarly, at $v=n$ we need $x_{n+1}=x_n$, and indeed $x_{n+1}=\cos(j\pi+\frac\theta2)$ and $x_n=\cos(j\pi-\frac\theta2)$ are both $(-1)^j\cos\frac\theta2$. Each $x$
is non-zero ($x_1=\cos\frac\theta2>0$), and the $n$ eigenvalues, for $j=0,\dots,n-1$, are distinct, so these are all the eigenvalues. The Fiedler vector ($j=1$) is a decreasing cosine; its sweep sets are the initial segments of the path, and the best is the
half path, cut in the middle.
\end{solution}

\begin{solution}{spectral:ex:moorebound}
\idea Grow balls around the edge: each new layer is at most $d-1$ times the previous one.
\writeup (a) Let $V_i$ be the set of vertices at distance $i$ from the edge $\{v_1,v_2\}$. As in the proof of Theorem~\ref{spectral:thm:alonboppana}, $|V_0|=2$ and $|V_{i+1}|\le(d-1)|V_i|$, so $|V_i|\le2(d-1)^i$, and summing over
$i\le r$ gives the bound. If the girth is at least $2r+2$, the vertices within distance $r$ of $v_1$ on the side away from $v_2$, and those of $v_2$, form two disjoint trees in which every vertex at
distance $i<r$ has $d-1$ children (a coincidence or an extra edge would close a cycle of length at most $2r+1$), so all inequalities are equalities.
(b) By (a) with $r=2k+2$, some vertex $x$ is at distance at least $2k+3$ from $\{v_1,v_2\}$. Take any edge $\{x,x'\}$: $x'$ is at distance at least $2k+2$. So the two edges are at distance at least
$2k+2$, and $\lambda_2\ge2\sqrt{d-1}-\frac{2\sqrt{d-1}-1}{k+1}$. For fixed $d$ and $n\to\infty$, $k$ can be taken of order $\log_{d-1}n$, which proves $\lambda_2\ge2\sqrt{d-1}-O(\frac1{\log n})$.
\end{solution}

\section*{Sources and further reading}

\begin{itemize}
\item \emph{The stories.} M.~S.~Pinsker, \emph{On the complexity of a concentrator}, 7th International Teletraffic Congress, Stockholm, 1973; G.~A.~Margulis, \emph{Explicit constructions of
concentrators}, Problems of Information Transmission 9 (1973), 325--332; O.~Reingold, \emph{Undirected connectivity in log-space}, Journal of the ACM 55 (2008), article 17; M.~Kac, \emph{Can one hear the shape of a drum?}, American Mathematical Monthly 73 (1966), 1--23.
\item G.~Kirchhoff, \emph{\"Uber die Aufl\"osung der Gleichungen, auf welche man bei der Untersuchung der linearen Vertheilung galvanischer Str\"ome gef\"uhrt wird}, Annalen der Physik
und Chemie 72 (1847), 497--508; W.~T.~Tutte, \emph{The dissection of equilateral triangles into equilateral triangles}, Proceedings of the Cambridge Philosophical Society 44 (1948),
463--482, for the directed theorem; T.~van Aardenne-Ehrenfest and N.~G.~de Bruijn, \emph{Circuits and trees in oriented linear graphs}, Simon Stevin 28 (1951), 203--217, for BEST.
\item M.~Fiedler, \emph{Algebraic connectivity of graphs}, Czechoslovak Mathematical Journal 23 (1973), 298--305; N.~Alon and V.~D.~Milman, \emph{$\lambda_1$, isoperimetric inequalities for graphs, and
superconcentrators}, Journal of Combinatorial Theory B 38 (1985), 73--88; J.~Dodziuk, \emph{Difference equations, isoperimetric inequality and transience of certain random walks}, Transactions
of the AMS 284 (1984), 787--794; J.~Shi and J.~Malik, \emph{Normalized cuts and image segmentation}, IEEE Transactions on Pattern Analysis and Machine Intelligence 22 (2000), 888--905.
\item A.~Nilli, \emph{On the second eigenvalue of a graph}, Discrete Mathematics 91 (1991), 207--210 (Nilli is a pseudonym of Noga Alon); N.~Alon and F.~R.~K.~Chung, \emph{Explicit construction
of linear sized tolerant networks}, Discrete Mathematics 72 (1988), 15--19; A.~Lubotzky, R.~Phillips and P.~Sarnak, \emph{Ramanujan graphs}, Combinatorica 8 (1988), 261--277; J.~Friedman,
\emph{A proof of Alon's second eigenvalue conjecture and related problems}, Memoirs of the AMS 195 (2008).
\item S.~Hoory, N.~Linial and A.~Wigderson, \emph{Expander graphs and their applications}, Bulletin of the AMS 43 (2006), 439--561: the survey to read next, written for a wide audience. A.~Marcus, D.~A.~Spielman
and N.~Srivastava, \emph{Interlacing families I: bipartite Ramanujan graphs of all degrees}, Annals of Mathematics 182 (2015), 307--325: a research paper of a dozen pages whose method, a careful use of
real-rooted polynomials and interlacing (Tome~2, Appendix~E), is within reach after this chapter.
\item D.~A.~Spielman, \emph{Spectral and Algebraic Graph Theory}, book draft, Yale, 2019 (freely available), the reference for this chapter's point of view; A.~E.~Brouwer and W.~H.~Haemers,
\emph{Spectra of Graphs}, Springer, 2012; O.~Bernardi, \emph{On the spanning trees of the hypercube and other products of graphs}, Electronic Journal of Combinatorics 19(4) (2012), P51.
\end{itemize}
